% Instrucciones incluidas en el manual de seguridad en radiología

\section*{INSTRUCCIONES DE OPERACIÓN}

Para unir el ducto de la fuente al cabezal, se quita el tapón con rosca (a en la figura 1) y en su lugar se atornilla el acople del extremo correspondiente del ducto (el cabezal debe estar bloqueado).

Para conectar el manipulador con indicador, se quita el tapón con rosca (b en la figura 1) y se conecta el cable de control al cable de la fuente, teniendo mucho cuidado en esta operación para que no se vaya a safar.

Para disponer el equipo para operación, el extremo del ducto de la fuente se coloca en la posición donde se va a tomar la radiografía, su fijación puede ser por un trípode, cinta adhesiva, abrazadera, bloque de madera con perforaciones u otra forma adecuada. El ducto debe estar tan recto como sea posible para una mejor facilidad de operación.

Una vez realizado lo anterior, se extiende el cable de control a toda su longitud y con la manija de control tan alejada como sea posible del extremo del ducto de la fuente en donde quedará expuesto el iridio. El operador también debe estar tan alejado como sea posible, del punto donde se va a utilizar la fuente, de preferencia atrás de un blindaje para el personal.

Después de abrir el iriditrón con la llave que se proporciona, se expone la fuente, girando la manija de control hasta que ésta se detenga. El indicador de posición muestra en todo momento, la distancia recorrida por la fuente. No se debe forzar la manija.

Una vez realizada la exposición, se regresa la fuente girando la manija en sentido contrario. Cierre la chapa de seguridad y retire la llave para que el equipo no sea utilizado por personal no autorizado.
